\documentclass{article}
\usepackage[T1]{fontenc}
\usepackage{lmodern}
\usepackage{graphicx}
\usepackage{amsmath,amssymb}
\usepackage[a4paper,margin=2cm]{geometry}
\usepackage[absolute,overlay]{textpos}
\setlength{\TPHorizModule}{1mm}
\setlength{\TPVertModule}{1mm}
\usepackage[indonesian]{babel}
\usepackage{hyperref}
\setlength{\emergencystretch}{2em}
% unit: Halaman 13

\begin{document}

% === Halaman 13 (llm) ===

\section*{Contoh pengiriman pesan 1:}

\begin{itemize}
    \item Misalkan Bob mengirim pesan $m = 16$ kepada Alice. Bob mengenkripsi $m$ dengan kunci publik Alice $(e, n) = (79, 3337)$ sebagai berikut:
    \[ c = m^e \text{ mod } n = 16^{79} \text{ mod } 3337 = 1493 \]
    
    Bob mengirim $c = 1493$ kepada Alice
    
    \item Alice mendekripsi cipherteks $c = 1493$ dengan kunci privat miliknya, yaitu $(d, n) = (1019, 3337)$ sebagai berikut:
    \[ m = c^d \text{ mod } n = 1493^{1019} \text{ mod } 3337 = 16 \]
    
    Alice mendapatkan kembali plainteks dari Bob, yaitu $m = 16$
\end{itemize}

\end{document}
